\summarypart{THREE}{The Age of the Christian Roman Empire}{312--590}
\begin{partoverview}
The emperor Constantine is one of the major figures of Christian history. After his conversion the
faith moved quickly from the seclusion of the catacombs to the prestige of palaces: it began the fourth
century as a persecuted minority and ended it as the established religion of the empire. Joined to the
power of the state, the church took on moral responsibility for the whole of society and, at first
under Constantine's direction, refined its doctrine and developed its structure.

Christians judged the change differently. Some, like the historian Eusebius, saw Christianity's victory
over the empire; others, like the monks, feared that the culture was capturing Christianity. What
follows is the ``Christianizing'' of the Hellenistic mind and world. And when the empire gave way to
the barbarian invaders---the ancestors of today's Europeans---it was, ironically, the monks who won
the conquerors, by the dignity of a well-ordered common life rooted in the Christian faith.
\end{partoverview}
\partchapters{%
  \pch{9}{Laying Her Sceptre Down}{The Conversion of the Empire}
  \pch{10}{Splitting Important Hairs}{The Doctrine of the Trinity}
  \pch{11}{Emmanuel!}{Christ in the Creeds}
  \pch{12}{Exiles from Life}{The Beginnings of Monasticism}
  \pch{13}{The Sage of the Ages}{Augustine}
  \pch{14}{Peter as ``Pontifex Maximus''}{The Beginnings of the Papacy}
  \pch{15}{Somewhere Between Heaven and Earth}{Eastern Orthodoxy}
  \pch{16}{Bending the Necks of Victors}{Mission to the Barbarians}}
